% from https://github.com/fwalch/tum-thesis-latex
\chapter{\abstractname}

Identifying which memory accesses limit a high-performance application's throughput requires precise, hardware-assisted sampling --- the role Intel PEBS and AMD IBS were designed to play. Mitos, originally developed at Lawrence Livermore National Laboratory as the data-collection front-end for the MemAxes visualiser, exposes these facilities to HPC developers: it samples loads above a configurable latency threshold, attributes each sample to a source file and line, and produces an output bundle that downstream tools consume. Its implementation, however, encodes PEBS events as literal opcodes against a specific Intel microarchitecture, which both ties the tool to a small set of supported CPUs and makes supporting new microarchitectures or sampling technologies a manual, error-prone effort.

This thesis re-implements Mitos on top of the Caliper performance instrumentation framework and the libpfm4 PMU-event library, replacing direct \texttt{perf\_event\_attr} construction in the tool itself with symbolic event encoding and the bespoke sampling machinery with a Caliper service, therefore cerating a small abstraction layer above archtecture-specific encodings. The same source compiles and runs unmodified on Intel Skylake and Intel Ice Lake; the elimination of the Dyninst dependency in favour of Caliper's symbol-lookup service and an output path compatible with the existing MemAxes pipeline follow as by-products of the redesign. An AMD platform was implemented, but access issues on the available systems prevented the test from completing in time; the IBS half of the cross-vendor claim therefore rests on the design argument alone.

Overall the new implementation is severely lacking compared to the original Mitos. It is briefly argued that the design is sound and the implementation is a first step towards a more robust, future-proof Mitos, but the current state of the codebase is not fit for production use. This argument is not detailed and unverified.
